Fix \(0<\Delta\leq\Delta_0\) and \(Q\in\mathcal Q_\Delta\).  By \(\mathrm{def:released\text{-}identified\text{-}set}\), there is at least one \(P\in\mathcal P_L\) such that \(Q=\mathcal L_P(Z_\Delta)\).  For \(t\in\{0,1\}\), define the released first-bin side event
\[
 B_t=\{T=t,Q_\Delta(R)=0\}.
\]
Because \(R\in[0,1]\) and \(Q_\Delta(r)=\Delta\lfloor r/\Delta\rfloor\),
\[
 B_t=\{T=t,0\leq R<\Delta\}=\{T=t,R<\Delta\}.
\]
The inequalities \(0<\Delta\leq\Delta_0\leq1/4<1/2\) permit use of the two-sided local-mass condition in \(\mathrm{ass:uniform\text{-}disc\text{-}design}\) and \(\mathrm{def:half\text{-}disc\text{-}class}\), giving
\[
 Q(B_t)=\mathbb P_P(B_t)\geq c_0\Delta^2>0.
\]
This proves the first assertion for every attainable \(Q\).

Since \(B_t\) has positive \(Q\)-probability, the released law determines the number
\[
 \zeta_t(Q)=\mathbb E_Q[Y\mid B_t]
 =\frac{\int y\mathbf1_{B_t}(y,t',q)\,dQ(y,t',q)}{Q(B_t)}.
\]
For any \(P'\in\mathcal P_L\) satisfying \(\mathcal L_{P'}(Z_\Delta)=Q\), the general event assertion of \(\mathrm{lem:conditional\text{-}local\text{-}binomial\text{-}law}\), which follows from sharp assignment, consistency, and the declared conditional means, yields
\[
 \zeta_t(Q)=\mathbb E_{P'}[Y\mid B_t]
 =\mathbb E_{P'}[\mu_t(X)\mid B_t].
\]
On \(B_t\), one has \(R<\Delta\leq1/2\).  The point-anchor envelope in \(\mathrm{ass:radial\text{-}lipschitz}\) therefore gives, using conditional Jensen's inequality for the absolute value,
\[
 \begin{aligned}
 |\zeta_t(Q)-\theta_{t,P'}|
 &\leq \mathbb E_{P'}[|\mu_t(X)-\theta_{t,P'}|\mid B_t]\\
 &\leq L\mathbb E_{P'}[R\mid B_t]
 \leq L\Delta.
 \end{aligned}
\]
The anchor \(\theta_{t,P'}\) is unambiguous by \(\mathrm{lem:boundary\text{-}anchor\text{-}uniqueness}\).  Hence, if \(P',P''\in\mathcal P_L\) induce the same \(Q\), then
\[
 |\theta_{t,P'}-\theta_{t,P''}|
 \leq |\theta_{t,P'}-\zeta_t(Q)|+|\zeta_t(Q)-\theta_{t,P''}|
 \leq2L\Delta.
\]
Using \(\tau_P(b)=\theta_{1,P}-\theta_{0,P}\) and the triangle inequality,
\[
 |\tau_{P'}(b)-\tau_{P''}(b)|
 \leq |\theta_{1,P'}-\theta_{1,P''}|+|\theta_{0,P'}-\theta_{0,P''}|
 \leq4L\Delta.
\]
Taking the supremum first over all \(P',P''\) in a fixed release fiber and then over \(Q\in\mathcal Q_\Delta\), as prescribed by \(\mathrm{def:released\text{-}identified\text{-}set}\), proves \(\omega_\Delta\leq4L\Delta\).

For the converse, \(\mathrm{lem:first\text{-}bin\text{-}alias}\) supplies \(P^{\mathrm{alias}}_{+,\Delta},P^{\mathrm{alias}}_{-,\Delta}\in\mathcal P_L\) having a common released law, say \(Q^{\mathrm{alias}}\), and satisfying
\[
 \tau_{P^{\mathrm{alias}}_{+,\Delta}}(b)
 -\tau_{P^{\mathrm{alias}}_{-,\Delta}}(b)=2c_L\Delta.
\]
Both displayed target values consequently belong to \(I_\Delta(Q^{\mathrm{alias}})\), so
\[
 \omega_\Delta
 \geq\operatorname{diam}I_\Delta(Q^{\mathrm{alias}})
 \geq2c_L\Delta.
\]
Here \(c_L>0\) by its definition and \(L\geq1\).  Thus a fixed-width release has at least one fiber containing two distinct values of the structural causal target, which proves general failure of point identification.  The two positive finite constants \(2c_L\) and \(4L\) give the asserted order in \(\Delta\).